% =====================================================================
% PDM — Patch-Diffusion Pipeline: Implementation-Grounded Walkthrough
% A reverse-engineering / teaching document built from the source code
% and the training / calibration / evaluation logs in outputs/.
%
% Compile:  pdflatex pipeline_walkthrough.tex   (run twice for the ToC)
%        or upload to Overleaf.  Standard packages only — no figures
%        are \include-d, so it compiles stand-alone.
% =====================================================================
\documentclass[11pt,a4paper]{article}

\usepackage[utf8]{inputenc}
\usepackage[a4paper,margin=22mm]{geometry}
\usepackage{amsmath,amssymb,amsthm,mathtools}
\usepackage{bm}
\usepackage{booktabs}
\usepackage{array}
\usepackage{longtable}
\usepackage{xcolor}
\usepackage{listings}
\usepackage{enumitem}
\usepackage{titlesec}
\usepackage[colorlinks=true,linkcolor=blue!50!black,urlcolor=red!50!black]{hyperref}

\newcommand{\R}{\mathbb{R}}
\newcommand{\E}{\mathbb{E}}
\newcommand{\N}{\mathcal{N}}
\newcommand{\x}{\bm{x}}
\newcommand{\eps}{\bm{\epsilon}}
\newcommand{\II}{\mathbf{I}}
\newcommand{\abar}{\bar{\alpha}}
\newcommand{\code}[1]{\texttt{\small #1}}

\definecolor{codebg}{gray}{0.96}
\lstset{
  basicstyle=\ttfamily\footnotesize,
  backgroundcolor=\color{codebg},
  breaklines=true, frame=single, framesep=4pt,
  columns=fullflexible, keepspaces=true, showstringspaces=false,
}

\title{\textbf{The Patch-Diffusion (PDM) Pipeline:\\
A Sequential, Implementation-Grounded Walkthrough}\\[4pt]
\large From a raw BraTS-PEDs volume to the final anomaly metrics}
\author{Reverse-engineered from the source code and the \code{outputs/} logs}
\date{\today}

\begin{document}
\maketitle

\begin{abstract}
\noindent
This document is a complete, sequential explanation of the pixel-space
\emph{patch-diffusion} pipeline for unsupervised anomaly detection (UAD) in
pediatric brain MRI, as it is \emph{actually implemented} in this repository. It
follows the data from a raw 3-D NIfTI volume, through preprocessing, patch
extraction, the forward diffusion process, the U-Net denoiser, training,
partial-noise reconstruction at inference, multi-scale residual scoring,
threshold calibration, evaluation, and finally explainability. Every claim is
tied to a specific file, function, configuration value, or log line. Where the
code and the logs are mutually inconsistent (there are two such places), the
ambiguity is stated explicitly rather than papered over. Dimensions are traced
layer-by-layer; the relevant mathematics is given with every symbol defined and
an intuitive reading for this specific project.
\end{abstract}

\tableofcontents
\clearpage

% =====================================================================
\section{System overview}
% =====================================================================

\subsection{What the system does, in one paragraph}
A denoising diffusion model is trained \emph{only} on $96\times96$ patches of
\emph{healthy} brain tissue, so it learns what healthy anatomy looks like. At
test time, each slice of an unseen patient is cut into overlapping patches; every
patch is partially corrupted with noise and then denoised back by the model. The
model can only reconstruct toward healthy anatomy, so healthy tissue returns
almost unchanged while a lesion is ``healed'' into plausible healthy tissue. The
pixel-wise difference between the original patch and its healthy reconstruction
(the \emph{residual}) is therefore large exactly where the anatomy is abnormal.
These residuals are fused back into a full-slice anomaly map, thresholded at a
value calibrated on healthy data, and compared against the radiologist's lesion
mask to produce AUROC, AUPRC and Dice.

\subsection{End-to-end stage map}
The pipeline is a linear sequence of stages, each owned by one module:

\begin{longtable}{@{}rll@{}}
\toprule
\# & Stage & Implementation \\
\midrule
\endhead
1  & Preprocess NIfTI $\to$ normalized 2-D slices & \code{src/data/preprocessing.py}, \code{normalization.py} \\
2  & Build healthy / val / test manifests        & \code{preprocessing.py} (writes \code{manifests/*.txt}) \\
3  & Patch index + dataset                        & \code{src/data/patches.py}, \code{dataset.py} \\
4  & Forward (noising) process                    & \code{src/models/diffusion.py}, \code{src/noise/simplex.py} \\
5  & U-Net $\eps$-predictor                        & \code{src/models/unet.py} (\code{diffusers.UNet2DModel}) \\
6  & Training loop (AdamW + EMA + AMP)            & \code{src/training/trainer.py}, \code{models/ema.py} \\
7  & Partial-noise DDIM reconstruction (inference)& \code{diffusion.py::reconstruct} \\
8  & Residual + modality weighting + CE channel   & \code{src/scoring/residual.py} \\
9  & Multi-scale scoring + Gaussian fusion        & \code{src/scoring/multiscale.py}, \code{data/patches.py} \\
10 & Threshold calibration on healthy slices      & \code{src/calibration/calibrator.py} \\
11 & Evaluation + metrics                         & \code{src/evaluation/evaluator.py}, \code{metrics.py} \\
12 & Explainability (counterfactual + attribution)& \code{src/xai/counterfactual.py}, \code{attribution.py} \\
\bottomrule
\end{longtable}

\noindent
All constants live in a single frozen-dataclass config,
\code{src/config.py}, exported as the shared instance \code{CONFIG}. Every value
quoted in this document is the default in that file unless a log shows it was
overridden at run time.

\subsection{The hardware and the two training runs (from the logs)}
The logs in \code{outputs/logs/} record \emph{two} runs on the same machine,
\textbf{NVIDIA A100-SXM4-80GB} (Google Colab, High-RAM):
\begin{itemize}[leftmargin=1.4em,itemsep=1pt]
\item \textbf{Pilot run} (\code{pilot\_training.log}, 2026-06-19): the full,
uncapped patch set --- \emph{2{,}088{,}400 train patches}, 5438 steps/epoch,
scheduled for 250 epochs.
\item \textbf{Final run} (\code{train.log}, 2026-06-19/20): a capped run ---
\emph{250{,}000 train patches/epoch}, 714 steps/epoch, 60 epochs, resumed from
epoch 8. This is the run whose checkpoint produced all reported results.
\end{itemize}

% =====================================================================
\section{Stage 1 — Data loading and preprocessing}
% =====================================================================

\subsection{Input: the four BraTS-PEDs modalities}
Each patient folder holds four co-registered, skull-stripped 3-D MRI volumes plus
an expert segmentation (\code{preprocessing.py::\_load\_patient}):
\code{t1n} (native T1), \code{t1c} (contrast-enhanced T1), \code{t2w} (T2),
\code{t2f} (T2-FLAIR), and \code{seg} (lesion labels, used \emph{only} for
evaluation). The channel order is fixed by
\code{CONFIG.data.modalities = (t1n, t1c, t2w, t2f)} and is relied upon
downstream (the contrast-enhancement channel is computed as channel-1 minus
channel-0).

\subsection{Robust intensity normalization}
Each modality volume is normalized independently
(\code{normalization.py::robust\_normalize\_volume}). Foreground (non-zero)
voxels only are used to compute the $[0.5, 99.5]$ percentiles
$(\ell, h)$; the volume is clipped to $[\ell,h]$ and mapped affinely to $[-1,1]$:
\begin{equation}
\tilde v \;=\; 2\,\frac{\mathrm{clip}(v,\ell,h)-\ell}{h-\ell}\;-\;1,
\qquad \tilde v\big|_{v=0} := -1 .
\end{equation}
\textbf{Symbols.} $v$ is a raw voxel intensity; $\ell,h$ are the low/high
foreground percentiles; $\tilde v\in[-1,1]$ is the normalized value; background
voxels (exactly $0$) are forced to the constant $-1$.

\textbf{Why this and not $\pm3\sigma$.} Percentile clipping at $99.5\%$ keeps the
hyper-intense FLAIR/enhancement signal that marks tumour. Hard $\sigma$-clipping
would truncate exactly that signal. The constant $-1$ background gives the model
a clean, predictable ``empty'' value, which later defines the brain mask.

\subsection{Slicing, resizing, and the foreground gate}
Volumes are sliced axially; each 2-D slice is centre-cropped/padded to
$256\times256$ with $-1$ padding (\code{\_resize\_slice}). The four modalities are
stacked into a $(4,256,256)$ \code{float16} array and saved as
\code{slices/<pid>\_z<NNN>.npy}. A slice is discarded if its foreground fraction
is below \code{min\_foreground\_frac}; note the config default is $0.02$ while the
docstring mentions a stricter intent --- the \emph{code} value $0.02$ is what
ran. Lesion slices additionally get a binary mask saved under \code{masks/}.

\subsection{Healthy filtering with a lesion buffer (leakage control \#1)}
A slice is admitted to the \emph{healthy} training pool only if it is lesion-free
\emph{and} at least \code{healthy\_buffer = 3} slices away from any lesion slice
(\code{\_healthy\_slice\_flags}). This buffer prevents tumour mass-effect and
partial-volume tissue immediately adjacent to a lesion from leaking into the
``healthy'' manifold; without it the model would partly learn peri-tumoural
appearance and then fail to flag it.

\subsection{Patient-level splits and the three manifests (leakage control \#2)}
Splits are read from \code{splits/\{train,val,test\}.txt} at the
\emph{patient} level, so no slice of a test patient is ever seen in training or
calibration. Preprocessing emits three manifests:
\begin{itemize}[leftmargin=1.4em,itemsep=0pt]
\item \code{healthy.txt} --- healthy+buffer slices from \emph{train} patients (diffusion training);
\item \code{val\_healthy.txt} --- lesion-free slices from \emph{val} patients (threshold calibration);
\item \code{test\_anom.txt} --- lesion slices from \emph{test} patients (evaluation).
\end{itemize}
The diffusion model and the residual score use \emph{no} masks at all; masks
enter only in the final metric computation. This is the structural reason the
reported numbers are not optimistically biased.

% =====================================================================
\section{Stage 2 — Patches and datasets}
% =====================================================================

\subsection{Why patches}
Whole-slice reconstruction UAD is dominated by brain-boundary edge artefacts and
is starved of data when only $\sim$30 patients are available. Operating on small
overlapping patches (a) multiplies the effective dataset size by two orders of
magnitude and (b) forces a \emph{local} texture model that cannot memorise the
global slice layout (\code{patches.py} docstring; Behrendt et al., 2023).

\subsection{Patch tiling --- exact geometry}
\code{iter\_patch\_coords(image\_size, patch\_size, stride)} tiles the slice with
overlap, clamping the final row/column so the edge is always covered. For the
actual values $256$, $96$, $16$:
\[
\text{tops} = \{0,16,32,\dots,160\}\;(11\text{ positions}),
\qquad \#\text{patches/slice} = 11^2 = \boxed{121}.
\]
\textbf{Note (code vs.\ comment).} The \code{PatchConfig} docstring estimates
``$\sim$196 patches/slice''; the \emph{actual} tiling produces $121$. The $121$
figure is what the implementation computes and what the patch-batch timing below
is consistent with.

\subsection{The two datasets}
\code{HealthyPatchDataset} (\code{dataset.py}) builds a flat
\code{(stem, top, left)} index over all healthy slices, using memory-mapped reads
so only each candidate patch region is touched, and \emph{drops background-only
patches} via \code{patch\_foreground\_fraction} $\ge$
\code{min\_patch\_foreground}$=0.10$. \code{\_\_getitem\_\_} materialises a single
$(4,96,96)$ \code{float32} patch and applies training augmentation:
horizontal flip (valid under approximate bilateral symmetry, $p{=}0.5$) and a
small multiplicative intensity jitter ($\pm5\%$, $p{=}0.3$).

\code{AnomalousSliceDataset} returns whole $(4,256,256)$ test slices with their
$(256,256)$ lesion masks; patching is deferred to the evaluator so that Gaussian
fusion can run over the whole slice.

% =====================================================================
\section{Stage 3 — The forward (noising) diffusion process}
% =====================================================================

\subsection{Closed-form forward kernel}
With a variance schedule $\beta_1,\dots,\beta_T$, define
$\alpha_t = 1-\beta_t$ and $\abar_t=\prod_{s\le t}\alpha_s$. The forward process
admits the closed form used in \code{diffusion.py::add\_noise} (which calls the
\code{diffusers} \code{DDPMScheduler.add\_noise}):
\begin{equation}
\label{eq:forward}
\x_t = \sqrt{\abar_t}\,\x_0 + \sqrt{1-\abar_t}\,\eps,
\qquad
q(\x_t\mid\x_0)=\N\!\big(\sqrt{\abar_t}\,\x_0,\;(1-\abar_t)\II\big).
\end{equation}
\textbf{Symbols.} $\x_0$ is a clean healthy patch; $t\in\{0,\dots,T-1\}$ is the
diffusion timestep; $\abar_t\in(0,1]$ is the cumulative signal-retention factor
($\abar_t\!\to\!1$ at $t{=}0$, $\abar_t\!\to\!0$ at $t{=}T$); $\eps$ is the
injected noise; $\II$ is the identity covariance.
\textbf{Intuition.} A single step jumps directly to any noise level $t$: keep a
fraction $\sqrt{\abar_t}$ of the clean image and add $\sqrt{1-\abar_t}$ of noise.

\subsection{Schedule (cosine)}
\code{CONFIG.model}: $T=1000$, \code{beta\_schedule = "squaredcos\_cap\_v2"}
(the Nichol \& Dhariwal cosine schedule), $\beta\in[10^{-4},2\times10^{-2}]$,
\code{prediction\_type = "epsilon"}. The cosine schedule keeps more signal at
high $t$ than a linear one, which empirically improves sample quality.

\subsection{Simplex (correlated) noise --- the key UAD ingredient}
Instead of white Gaussian noise, the forward process uses spatially
\emph{correlated} simplex-like noise (\code{noise/simplex.py}). It is built as a
sum of \code{octaves}$=6$ band-limited octaves: low-resolution Gaussian noise
bicubically upsampled to $96\times96$, with amplitude decayed by
\code{persistence}$=0.8$ and frequency grown by \code{lacunarity}$=2.0$ per
octave, starting at \code{base\_frequency}$=16$. The sum is renormalized to unit
variance per sample:
\begin{equation}
\bm n = \sum_{o=0}^{5} p^{o}\,\mathrm{up}\!\big(\bm g_o\big),
\qquad
\eps = \bm n \big/ \mathrm{std}(\bm n),
\end{equation}
where $\bm g_o$ is low-res Gaussian noise at octave $o$, $p=0.8$, and
$\mathrm{up}(\cdot)$ is bicubic upsampling to patch size.
\textbf{Why.} Correlated noise lives at the \emph{spatial scale of lesions}, so a
lesion ``looks like noise'' to the denoiser and is in-painted with healthy
tissue at test time, producing a large residual --- whereas white noise mostly
attacks pixel-level detail. Unit-variance renormalization keeps the DDPM
signal-to-noise schedule in Eq.~\eqref{eq:forward} valid.

% =====================================================================
\section{Stage 4 — The U-Net $\eps$-predictor (layer-by-layer)}
% =====================================================================

\subsection{Configuration}
\code{unet.py::build\_unet} instantiates a \code{diffusers.UNet2DModel}. The saved
\code{checkpoints/unet/config.json} confirms the exact built network:

\begin{longtable}{@{}ll@{}}
\toprule
Field & Value \\
\midrule
\code{sample\_size}        & 96 \\
\code{in\_channels / out\_channels} & 4 / 4 \\
\code{block\_out\_channels}& $[64, 128, 256, 256]$ \\
\code{layers\_per\_block}  & 2 \\
\code{down\_block\_types}  & Down, Down, Down, \textbf{AttnDown} \\
\code{up\_block\_types}    & \textbf{AttnUp}, Up, Up, Up \\
\code{mid\_block\_type}    & UNetMidBlock2D (ResNet–Attn–ResNet) \\
\code{attention\_head\_dim}& 8 \\
\code{act\_fn}             & SiLU \\
\code{norm\_num\_groups} / \code{norm\_eps} & 32 / $10^{-5}$ \\
\code{time\_embedding\_type}& positional (sinusoidal) \\
\code{dropout}            & 0.1 \\
Parameter count (logged)  & \textbf{25.3\,M} \\
\bottomrule
\end{longtable}

\noindent
The level resolutions are $96\!\to\!48\!\to\!24\!\to\!12$. Because
\code{attention\_resolutions}$=(12,)$, self-attention is applied \emph{only} at
the $12\times12$ bottleneck level (and in the mid-block), keeping memory and
parameter count low --- the right trade-off when training data is scarce.

\subsection{What the network predicts}
Given a noised patch $\x_t$ and timestep $t$, the network outputs
$\eps_\theta(\x_t,t)\in\R^{4\times96\times96}$: an estimate of the noise that was
added. This is the \code{.sample} field used in
\code{diffusion.py::training\_loss} and \code{reconstruct}.

\subsection{Dimensional walkthrough}
For a batch of $B$ patches, tensors flow as follows. Downsampling happens
\emph{after} each of the first three levels; the last (attention) down-level does
not downsample, which is why the bottleneck is $12\times12$.

\begin{longtable}{@{}lll@{}}
\toprule
Stage & Output shape $(C,H,W)$ & Notes \\
\midrule
\endhead
Input patch                 & $4\times96\times96$   & 4 modalities \\
\code{conv\_in} ($3{\times}3$) & $64\times96\times96$ & lift to base width \\
Time embed                  & vector $256$          & sinusoidal($64$)$\to$MLP$\to$256, added in every ResBlock \\
DownBlock 0 ($2\times$ResNet)& $64\times96\times96$ & then downsample $\to 64\times48\times48$ \\
DownBlock 1 ($2\times$ResNet)& $128\times48\times48$& then downsample $\to 128\times24\times24$ \\
DownBlock 2 ($2\times$ResNet)& $256\times24\times24$& then downsample $\to 256\times12\times12$ \\
\textbf{AttnDown 3} ($2\times$ResNet+Attn)& $256\times12\times12$ & self-attention, no downsample \\
Mid (ResNet–Attn–ResNet)    & $256\times12\times12$ & global mixing at bottleneck \\
\textbf{AttnUp 3} (+skip, Attn)& $256\times24\times24$ & concat skip, attention, upsample \\
UpBlock 2 (+skip)           & $256\times48\times48$ & concat skip, upsample \\
UpBlock 1 (+skip)           & $128\times96\times96$ & concat skip, upsample \\
UpBlock 0 (+skip)           & $64\times96\times96$  & concat skip (no further up) \\
GroupNorm+SiLU+\code{conv\_out}& $4\times96\times96$ & predicted noise $\eps_\theta$ \\
\bottomrule
\end{longtable}

\textbf{Residual blocks.} Each ResNet block is
GroupNorm($32$)$\to$SiLU$\to$Conv$3{\times}3$, with the time embedding projected
and \emph{added} after the first convolution, then
GroupNorm$\to$SiLU$\to$Dropout($0.1$)$\to$Conv$3{\times}3$ and a skip addition
(with a $1{\times}1$ conv when channel counts change). This is the standard DDPM
residual block.
\textbf{Skip connections.} The encoder feature maps at $96,48,24,12$ are
concatenated onto the matching decoder level (doubling channels at the
concatenation, which the first conv of each up-block reduces again), preserving
fine spatial detail that pure up-convolution would lose.
\textbf{Attention.} At $12\times12$, multi-head self-attention ($8$ heads) lets
every spatial location attend to every other --- affordable because there are
only $144$ tokens.

% =====================================================================
\section{Stage 5 — Training}
% =====================================================================

\subsection{Objective}
\code{diffusion.py::training\_loss} samples a uniform timestep per patch,
noises it via Eq.~\eqref{eq:forward}, and minimises the simple DDPM
$\eps$-prediction loss:
\begin{equation}
\label{eq:loss}
\mathcal L_{\text{simple}}
= \E_{\x_0,\;t\sim\mathcal U\{0,T-1\},\;\eps}
\big\|\,\eps - \eps_\theta(\x_t,t)\,\big\|_2^2 .
\end{equation}
\textbf{Reading the value.} Because $\eps$ is renormalized to unit variance, a
loss of $\mathcal L$ means the model leaves a fraction $\mathcal L$ of the noise
variance unexplained; the best validation loss $0.0604$ therefore corresponds to
$\approx94\%$ of the injected-noise variance explained. (This is an intuitive
reading, not an identity --- it holds because the target has unit variance.)

\subsection{Optimisation ingredients (\code{trainer.py})}
\begin{longtable}{@{}ll@{}}
\toprule
Component & Setting \\
\midrule
Optimizer        & AdamW, $\beta=(0.9,0.99)$, weight decay $10^{-4}$ \\
Peak LR          & $2\times10^{-4}$ \\
Schedule         & Linear warmup ($1$ epoch, start factor $0.1$) $\to$ cosine to $10^{-6}$ \\
EMA              & decay $0.9999$ (shadow weights, used for validation \emph{and} inference) \\
Mixed precision  & bf16 autocast (no GradScaler --- bf16 has fp32 range) \\
Grad clipping    & global-norm $1.0$ \\
Early stopping   & patience $6$ on EMA validation loss \\
Checkpointing    & best-val + periodic every 10 ep + full crash-safe state every ep \\
\bottomrule
\end{longtable}

\textbf{EMA matters here.} Validation (\code{\_validate}) and all downstream
inference use the EMA shadow module, not the raw weights. EMA averages out the
optimisation noise of the last steps and gives markedly cleaner reconstructions
--- which directly improves residual quality.

\textbf{Crash-safe resume.} \code{\_save\_full\_state} snapshots model+EMA+optimizer
+scheduler+bookkeeping every epoch, so a Colab disconnect costs at most one
epoch. The final run's log confirms this: it ``\emph{Resumed from epoch 8}''.

\subsection{What the logs actually show (final 60-epoch run)}
\begin{itemize}[leftmargin=1.4em,itemsep=1pt]
\item Device: \textbf{A100-SXM4-80GB}; \emph{Train patches} $=250{,}000$,
\emph{Val patches} $=7744$.
\item \textbf{714 optimiser steps/epoch}, $\approx565$\,s/epoch
($\approx9.4$\,min), total $\approx9.4$\,h for 60 epochs.
\item Validation loss falls smoothly: $0.17$ (ep 9, resumed) $\to 0.094$ (ep 25)
$\to 0.074$ (ep 34) $\to 0.0625$ (ep 50) $\to$ \textbf{best $0.06040$ at epoch 58}.
\item Final line: \code{DONE — best val 0.06040 @ epoch 58}; epoch 60 val
$0.06093$.
\end{itemize}

\textbf{Convergence reading.} Training loss ($\approx0.042$) and validation loss
($\approx0.060$) track each other with a small, stable gap and no upward
divergence --- i.e.\ \emph{neither under- nor over-fitting}. The $\sim$570k-patch
amplification is doing its job: the model is data-rich at the patch level even
though it has seen few patients. This is important later: it lets the failure
analysis \emph{exclude} training pathology as the cause of weak segmentation.

\subsection{Ambiguity: the logged batch size is unreliable}
\code{trainer.py} logs \code{CONFIG.train.batch\_size} (``\code{bs 128}''), but the
\emph{DataLoader} in \code{scripts/01\_train.py} is built with \code{args.bs}
(which defaults to the config but can be overridden on the command line). The
step counts prove the printed value is not the effective one:
\[
\text{final run: } \frac{250{,}000}{714}\approx 350,
\qquad
\text{pilot run: } \frac{2{,}088{,}400}{5438}\approx 384 .
\]
So the \textbf{effective batch size was $\approx$350 (final) / $\approx$384
(pilot)}, not $128$. The ``\code{bs 128}'' line is a logging artefact (it prints
the config default rather than \code{args.bs}). Any hyperparameter table should
quote the effective value, or simply ``$\approx$350'', rather than the logged
128.

% =====================================================================
\section{Stage 6 — Inference: partial-noise DDIM reconstruction}
% =====================================================================

\subsection{The reconstruction primitive}
\code{diffusion.py::reconstruct} is the inference heart. A clean test patch
$\x_0$ is noised \emph{only} to an intermediate level $t_{\text{int}}$ (not to
pure noise), then deterministically denoised with DDIM ($\eta=0$):
\begin{enumerate}[leftmargin=1.5em,itemsep=1pt]
\item add noise to $t_{\text{int}}$ via the DDIM scheduler's \code{add\_noise}
(same closed form as Eq.~\eqref{eq:forward}, with \emph{simplex} noise);
\item iterate the DDIM update from $t_{\text{int}}$ down to $0$.
\end{enumerate}

\subsection{DDIM update}
At each step the model first forms the clean-image estimate
\begin{equation}
\hat\x_0(\x_t,t)=\frac{\x_t-\sqrt{1-\abar_t}\,\eps_\theta(\x_t,t)}{\sqrt{\abar_t}},
\end{equation}
then takes the deterministic ($\eta=0$) DDIM step
\begin{equation}
\x_{t-1}=\sqrt{\abar_{t-1}}\;\hat\x_0(\x_t,t)\;+\;\sqrt{1-\abar_{t-1}}\;\eps_\theta(\x_t,t).
\end{equation}
\textbf{Symbols.} $\hat\x_0$ is the running estimate of the healthy clean patch;
$\eps_\theta$ is the network's noise prediction; $\eta=0$ makes the trajectory
deterministic (no added randomness), so the only stochasticity is the initial
noise draw.
\textbf{Intuition.} Healthy texture is on the manifold the model learned, so it
is reconstructed almost exactly; a lesion is off-manifold and gets pulled toward
the nearest healthy appearance --- it is ``healed''.

\subsection{How many reverse steps actually run}
\code{\_inference\_timesteps} sets the full DDIM grid to
$\max\!\big(\text{ddim\_steps},\,\mathrm{round}(\text{ddim\_steps}\cdot T/t_{\text{int}})\big)$
and then keeps only timesteps $\le t_{\text{int}}$. With
\code{ddim\_steps}$=50$ and $T=1000$ this yields $\approx$\textbf{50 reverse
steps at every scale}:
\[
t_{\text{int}}{=}50\!:\,\text{grid }1000\!\to\!51\text{ kept};\quad
t_{\text{int}}{=}150\!:\,\text{grid }333\!\to\!51;\quad
t_{\text{int}}{=}250\!:\,\text{grid }200\!\to\!50 .
\]
This is a neat design: the step \emph{count} is held roughly constant across
scales while the step \emph{size} adapts, so each scale costs the same compute.

% =====================================================================
\section{Stage 7 — Residual, modality weighting, and the CE channel}
% =====================================================================

\subsection{Per-channel residual}
\code{residual.py::residual\_stack} computes the absolute reconstruction error
per modality and appends a contrast-enhancement (CE) residual channel:
\begin{equation}
r_k=\big|\x_0^{(k)}-\hat\x_0^{(k)}\big|\ (k=1..4),
\qquad
r_{\text{CE}}=\Big|\big(\x_0^{(2)}-\x_0^{(1)}\big)-\big(\hat\x_0^{(2)}-\hat\x_0^{(1)}\big)\Big|,
\end{equation}
where channels $1{=}t1n$, $2{=}t1c$. The CE term tracks errors in the
\emph{gadolinium enhancement} signal $t1c-t1n$, a primary marker of active
tumour that is invisible to pipelines dropping native T1. The saved
\code{numpy\_dumps/*.npz} confirm a $(5,256,256)$ residual stack (4 modalities +
CE) and a $(4,256,256)$ healthy reconstruction.

\subsection{Modality-weighted scalar map}
\code{weighted\_residual} normalizes the salience weights
$\bm w\propto(0.5,1.0,0.75,1.5)$ for $(t1n,t1c,t2w,t2f)$ plus
$w_{\text{CE}}=1.0$, then forms the brain-masked weighted sum:
\begin{equation}
S_{\text{patch}}(\bm u)=\Big(\sum_{k}\bar w_k\,r_k(\bm u)\Big)\,m_{\text{brain}}(\bm u),
\qquad \bar w_k=\frac{w_k}{\sum_j w_j}.
\end{equation}
FLAIR ($t2f$, weight $1.5$) and contrast ($t1c$, weight $1.0$) dominate, matching
where tumour signal lives. The brain mask
(\code{residual.py::brain\_mask\_2d}) is the foreground (any channel $>-1$),
hole-filled, then eroded by $2$\,px to drop the skull-strip rim while keeping
peripheral cortical lesions.

% =====================================================================
\section{Stage 8 — Multi-scale scoring and Gaussian fusion}
% =====================================================================

\subsection{Per slice}
\code{multiscale.py::score\_slice}:
\begin{enumerate}[leftmargin=1.5em,itemsep=1pt]
\item extract the $121$ overlapping patches of the slice;
\item for each scale $T\in\{50,150,250\}$, reconstruct every patch
(mini-batched, \code{patch\_batch}$=64$) and compute its weighted residual;
\item fuse the patch residuals into a full-slice map with a centred Gaussian
window;
\item combine the three single-scale maps with weights $(0.25,0.45,0.30)$.
\end{enumerate}

\subsection{Gaussian fusion}
\code{patches.py::GaussianPatchFuser} weights each patch's contribution by a
centred Gaussian ($\sigma=32$\,px) so a patch's confident centre dominates its
context-starved edges. The fused value at pixel $\bm u$ is the weight-normalized
average over all patches covering it:
\begin{equation}
S_T(\bm u)=\frac{\sum_i g(\bm u-\bm c_i)\,S^{(i)}_{\text{patch}}(\bm u)}
{\sum_i g(\bm u-\bm c_i)},
\qquad
g(\bm d)=\exp\!\Big(\!-\tfrac{\|\bm d\|^2}{2\sigma^2}\Big),
\end{equation}
with $\bm c_i$ the centre of patch $i$. This averaging is what removes the
boundary edge artefacts that plague whole-image UAD.

\subsection{Multi-scale combination}
\begin{equation}
\label{eq:multiscale}
S(\bm u)=\Big(\sum_{T}\omega_T\,S_T(\bm u)\Big)\,m_{\text{brain}}(\bm u),
\qquad \omega=(0.25,0.45,0.30).
\end{equation}
Small $T$ (high SNR) is sensitive to fine/texture anomalies; large $T$ (low SNR)
to large structural ones. The mid scale $T{=}150$ carries the most weight.

\subsection{Cost accounting (matches the eval timing)}
Per slice: $121$ patches $\times\,3$ scales $\times\,\approx50$ DDIM steps
$\approx 1.8\times10^4$ U-Net forward passes, mini-batched in groups of $64$.
The evaluation log shows $\approx$\textbf{92\,s/slice}, consistent with this
budget on the A100.

% =====================================================================
\section{Stage 9 — Threshold calibration}
% =====================================================================
\code{calibrator.py} runs the \emph{same} multi-scale scorer on
\code{val\_healthy} slices (which should score near zero), pools all healthy
brain-voxel scores, and sets the operating threshold to their $95$th percentile:
\begin{equation}
\tau=\mathrm{percentile}_{95}\big(\{S(\bm u): m_{\text{brain}}(\bm u)=1,\ \text{healthy}\}\big).
\end{equation}
The saved \code{outputs/calibration.json} records the realised values:
$\tau=\mathbf{0.1280}$, healthy-score mean $0.0535$, std $0.0385$, from
$n=200$ healthy slices. \textbf{Reading $\tau$.} A test voxel above $0.128$ is
more anomalous than $\approx95\%$ of healthy tissue, giving a controlled,
label-free false-positive rate. Calibrating on held-out healthy val patients
only (never test, never lesions) is leakage control \#4.

% =====================================================================
\section{Stage 10 — Evaluation and metrics}
% =====================================================================

\subsection{Per-slice procedure}
\code{evaluator.py::evaluate} loads $\tau$, then for each test slice: scores it,
optionally CRF-refines (the \emph{decision} thresholds the raw score, not the CRF
map), thresholds at $\tau$ to get the prediction, and computes AUROC, AUPRC, Dice,
plus an oracle Dice from a $60$-point threshold grid.

\subsection{Metric definitions (\code{metrics.py})}
\begin{align}
\text{Dice}(P,G) &= \frac{2|P\cap G|}{|P|+|G|},\\
\text{AUROC} &= \Pr\!\big(S(\bm u^+)>S(\bm u^-)\big)\ \text{over brain voxels},\\
\text{AUPRC} &= \text{average precision over brain voxels},\\
\text{oracle Dice} &= \max_{\theta\in\text{grid}}\ \text{Dice}\big(\mathbb 1[S>\theta],G\big).
\end{align}
\textbf{Symbols.} $P$ is the predicted lesion set, $G$ the ground-truth lesion
set; $\bm u^+/\bm u^-$ are lesion/healthy voxels. AUROC/AUPRC are
\emph{threshold-free} (they judge ranking quality); Dice is at the deployable
$\tau$; oracle Dice is the per-slice best-threshold ceiling. The gap
calibrated$\to$oracle quantifies how much is left on the table by a single global
threshold.

\subsection{What the evaluation actually covered}
The eval log is indexed \code{[\,k/1000\,]} --- a $1000$-slice run was
\emph{planned} --- but the last logged line is \code{[205/1000]}. Evaluation was
therefore \textbf{truncated at $\approx$205 slices} (204 with valid metrics),
spanning four test patients in order: \texttt{P00203, P00002, P00216, P00076}.
This is precisely why every result is reported over \textbf{204 slices / 4
patients} rather than the planned 1000. The truncation is the single biggest
caveat on the headline numbers (see \S\ref{sec:reading}).

\subsection{Aggregate and per-patient results}
\begin{longtable}{@{}lrcccc@{}}
\toprule
Patient & Slices & AUROC & AUPRC & Dice (calib.) & Oracle Dice \\
\midrule
\endhead
P00203 & 61 & 0.750 & 0.047 & 0.053 & 0.108 \\
\textbf{P00002} & 36 & \textbf{0.893} & \textbf{0.222} & \textbf{0.270} & \textbf{0.321} \\
P00216 & 53 & 0.571 & 0.039 & 0.035 & 0.073 \\
P00076 & 54 & 0.692 & 0.046 & 0.013 & 0.112 \\
\midrule
\textbf{Overall} & \textbf{204} & \textbf{0.713} & \textbf{0.076} & \textbf{0.076} & \textbf{0.138} \\
\bottomrule
\end{longtable}
\noindent
The single best slice, \texttt{P00002\_z094}, reaches AUROC $0.965$ / Dice
$0.589$ (saved in \code{numpy\_dumps/}). The worst diagnostic slice,
\texttt{P00216\_z037}, has AUROC $0.44$.

\subsection{Interpreting the numbers}
\label{sec:reading}
\begin{itemize}[leftmargin=1.4em,itemsep=1pt]
\item \textbf{Score separability, not thresholding, is the bottleneck.} The
calibrated$\to$oracle Dice gap ($0.076\!\to\!0.138$) is only $\approx0.06$; on
P00216 even the oracle Dice is $0.073$ because many slices have AUROC $<0.5$ ---
no threshold can rank a lesion that the score does not separate.
\item \textbf{Strong bimodality.} P00002 (AUROC $0.893$) vs.\ P00216 (AUROC
$0.571$). With only four test patients, one patient moves the mean substantially;
the aggregate has no confidence interval and is not statistically robust.
\item \textbf{Likely root cause (two diagnostic cases).} On the failure slice the
residual is high across the \emph{whole} parenchyma and diffuse across
\emph{every} modality (\code{numpy\_dumps} + \code{report\_figures}), i.e.\ the
model cannot reconstruct that patient's healthy tissue well --- an out-of-manifold
anatomy effect expected from a $30$-patient training set. This is a hypothesis
from two cases, not a controlled result.
\end{itemize}

% =====================================================================
\section{Stage 11 — Explainability}
% =====================================================================

\subsection{Counterfactual (primary)}
\code{xai/counterfactual.py}: the healthy reconstruction $\hat\x$ \emph{is} the
decision variable, so the anomaly map $|\x-\hat\x|$ is faithful by construction
(no post-hoc surrogate). $\hat\x$ is the projection of the patient onto the
learned healthy manifold --- ``what this brain would look like without disease''.
The module also saves a DDIM ``healing trajectory'' (\code{trajectory\_frames}$=6$)
showing the lesion being progressively in-painted.

\subsection{Exact additive attribution (supporting)}
Because the score in Eq.~\eqref{eq:multiscale} is an exact sum over modalities
and scales, attribution is exact, not approximate:
\begin{itemize}[leftmargin=1.4em,itemsep=0pt]
\item \code{attribution.py::modality\_attribution} splits the score into the five
weighted channel maps ($t1n,t1c,t2w,t2f,$CE) that sum back to it --- answering
``is this lesion mainly FLAIR or enhancement?'';
\item \code{scale\_attribution} splits it by noise scale $T$ --- low-$T$ dominance
$\Rightarrow$ fine anomaly, high-$T$ $\Rightarrow$ structural.
\end{itemize}
On the success slice FLAIR dominates the genuine lesion; on the failure slice the
signal spreads evenly across modalities --- the signature of a whole-brain
reconstruction error rather than focal pathology.

\subsection{Attention rollout (wired, not implemented)}
\code{attribution.py::AttentionRollout} and the \code{attention\_rollout} config
flag exist, but the class captures channel-mean feature magnitudes via forward
hooks (a best-effort proxy) and is not part of the reported results. Tracing the
true $12\times12$ self-attention back to patch positions is left as future work.

% =====================================================================
\section{End-to-end: one slice from raw input to metric}
% =====================================================================
Putting every stage together for a single test slice $\x_{\text{slice}}\in\R^{4\times256\times256}$:
\begin{enumerate}[leftmargin=1.6em,itemsep=2pt]
\item \textbf{Mask}: $m_{\text{brain}}=$ hole-filled, $2$-px-eroded foreground.
\item \textbf{Tile}: extract $121$ patches $\x\in\R^{4\times96\times96}$ (stride 16).
\item \textbf{For each scale} $T\in\{50,150,250\}$:
  \begin{enumerate}[leftmargin=1.3em,itemsep=0pt]
  \item noise each patch to $T$ with simplex noise; DDIM-denoise ($\eta{=}0$,
  $\approx50$ steps) with the EMA U-Net $\to\hat\x_0$;
  \item residual stack $r_k$ + CE; modality-weighted scalar $S_{\text{patch}}$;
  \item Gaussian-fuse the 121 patch maps $\to S_T$ (full slice).
  \end{enumerate}
\item \textbf{Combine scales}: $S=\sum_T\omega_T S_T$, masked.
\item \textbf{Threshold}: $\hat P=\mathbb 1[S>\tau]\cdot m_{\text{brain}}$,
$\tau=0.128$.
\item \textbf{Score}: AUROC/AUPRC from $S$ vs.\ $G$; Dice from $\hat P$ vs.\ $G$;
oracle Dice from the threshold grid.
\item \textbf{Explain}: counterfactual $\hat\x$, $|\x-\hat\x|$, and the
per-modality / per-scale attribution maps.
\end{enumerate}
The model is trained \emph{once} on healthy patches (Stage 5); steps 1--7 use only
the frozen EMA weights and the single calibrated scalar $\tau$. No lesion label
ever touches the model --- masks appear only in step 6.

% =====================================================================
\appendix
\section{Source-file map}
\begin{longtable}{@{}ll@{}}
\toprule
File & Contributes \\
\midrule
\code{src/config.py} & All constants (frozen dataclasses); shared \code{CONFIG}. \\
\code{data/preprocessing.py} & NIfTI$\to$slices, healthy buffer, manifests, splits. \\
\code{data/normalization.py} & Robust $[0.5,99.5]$ percentile $\to[-1,1]$. \\
\code{data/patches.py} & Patch tiling + \code{GaussianPatchFuser}. \\
\code{data/dataset.py} & \code{HealthyPatchDataset}, \code{AnomalousSliceDataset}. \\
\code{noise/simplex.py} & Multi-octave correlated noise (unit-variance). \\
\code{models/unet.py} & Builds \code{diffusers.UNet2DModel} (25.3\,M params). \\
\code{models/diffusion.py} & Forward kernel, loss, DDIM \code{reconstruct}. \\
\code{models/ema.py} & EMA shadow weights (used at val + inference). \\
\code{training/trainer.py} & AdamW+warmup/cosine+EMA+AMP loop, resume, CSV. \\
\code{scoring/residual.py} & Brain mask, modality weights, CE channel. \\
\code{scoring/multiscale.py} & Per-scale reconstruct+fuse $\to$ \code{ScoreResult}. \\
\code{calibration/calibrator.py} & $95$th-percentile healthy threshold ($\tau$). \\
\code{evaluation/evaluator.py} & Per-slice scoring, thresholding, panels. \\
\code{evaluation/metrics.py} & Dice, AUROC, AUPRC, oracle grid. \\
\code{xai/counterfactual.py} & Healthy counterfactual + healing trajectory. \\
\code{xai/attribution.py} & Exact modality/scale attribution; attention rollout. \\
\bottomrule
\end{longtable}

\section{Key realised numbers (from \code{outputs/})}
\begin{longtable}{@{}ll@{}}
\toprule
Quantity & Value (source) \\
\midrule
GPU & A100-SXM4-80GB (\code{train.log}) \\
Train patches / epoch & 250{,}000 (final run); 2{,}088{,}400 (pilot) \\
Val patches & 7{,}744 \\
Steps / epoch & 714 (final); 5{,}438 (pilot) \\
Effective batch size & $\approx$350 (final); $\approx$384 (pilot) --- \emph{not} the logged 128 \\
Time / epoch & $\approx$565\,s ($\approx$9.4\,min) \\
Epochs & 60 (resumed from 8) \\
Best val loss & \textbf{0.06040 at epoch 58} \\
Calibrated $\tau$ & 0.1280 ($n=200$ healthy slices) \\
Healthy score mean / std & 0.0535 / 0.0385 \\
Eval coverage & 204 valid slices / 4 patients (planned 1000, truncated at $\sim$205) \\
Time / slice (eval) & $\approx$92\,s \\
Overall AUROC / AUPRC / Dice / oracle & 0.713 / 0.076 / 0.076 / 0.138 \\
Best slice (P00002\_z094) & AUROC 0.965 / Dice 0.589 \\
\bottomrule
\end{longtable}

\end{document}
